\begin{khmer}
ការរាប់គ្រឿងបង្គុំខ្នាត​​តូចប្រកបដោយភាពជាក់លាក់ និងទាន់ពេលវេលា នៅលើខ្សែបញ្ជូនក្នុងឧស្សាហកម្មទំនើបនិងការវិវត្ត​ទៅរកបច្ចេកវិទ្យាបញ្ញាសប្បនិមិត្ត នៅតែជាបញ្ហាប្រឈមយ៉ាងខ្លាំងនៅក្នុងលូជីស្ទីកផលិតកម្ម។ ការត្រួតពិនិត្យដោយការមើលផ្ទាល់គឺតម្រូវឱ្យមានការប្រើប្រាស់ខួរក្បាល និងការផ្តោតអារម្មណ៍ខ្ពស់ ហើយងាយមានកំហុស ខណៈដែលដំណោះស្រាយស្វ័យប្រវត្តិកម្មបែបប្រពៃណីតែងតែជួបបញ្ហាដោយសារពន្លឺចាំងផ្លាស់ ការត្រួតស៊ីគ្នាយ៉ាងក្រាស់នៃវត្ថុ និងការចំណាយធនធានគណនាដ៏ច្រើនហួសហេតុរបស់ម៉ូដែលចាប់សញ្ញាដែលមានភាពជាក់លាក់ខ្ពស់។ បញ្ហាប្រឈមទាំងនេះជាសរុប បានរារាំងការដាក់ឱ្យដំណើរការប្រកបដោយភាពជឿជាក់ និងជាប់លាប់ នៅលើប្រព័ន្ធ \textenglish{Edge} របស់រោងចក្រដែលមានកម្រិតធនធាន។\\

ដើម្បីដោះស្រាយបញ្ហានេះ និក្ខេបបទនេះបានស្នើឡើងនូវក្របខ័ណ្ឌ \textenglish{Edge-AI} ទម្ងន់ស្រាលមួយ សម្រាប់ការចាប់សញ្ញា និងរាប់គ្រឿងចាប់ហាប់លោហៈទាន់ពេលវេលា។ សំណុំទិន្នន័យមួយត្រូវបានបង្កើតឡើងជាពិសេសនៅក្នុងបរិស្ថានម៉ូដែលរូបវន្ត ដែលក្លែងធ្វើតាមលក្ខខណ្ឌនៃខ្សែបញ្ជូនគោលដៅ ដោយរួមបញ្ចូលទាំងការប្រែប្រួលនៃពន្លឺ ភាពព្រាលដោយសារចលនា និងវត្ថុដែលកកកុញយ៉ាងណែន។ ដើម្បីបង្កើតមូលដ្ឋានគោលយ៉ាងតឹងរ៉ឹង និងកំណត់ការរឹតត្បិតនៃស្ថាបត្យកម្មនៅក្នុងការចាប់សញ្ញាវត្ថុខ្នាតតូច ការកំណត់រចនាសម្ព័ន្ធនៃម៉ូដែលចាប់សញ្ញាចំនួនបួនត្រូវបានធ្វើតេស្តប្រៀបធៀប៖ \textenglish{YOLO11s, YOLO11m, YOLOv26m}, និង \textenglish{RT-DETR} ជាឯកសារយោងដែលមានភាពជាក់លាក់ខ្ពស់។\\

ដោយផ្អែកលើការរឹតត្បិតនៃមូលដ្ឋានគោល ការងារនេះបានស្នើឡើងនូវ \textenglish{YOLOv8-GED} ដែលជាម៉ូដែលចាប់សញ្ញាទម្ងន់ស្រាលដែលត្រូវបានរចនាឡើងជាពិសេស និងប្រើប្រាស់ឱ្យមានប្រសិទ្ធភាពបំផុតសម្រាប់ការដាក់ឱ្យដំណើរការលើប្រព័ន្ធ \textenglish{Edge}។ ស្ថាបត្យកម្មនេះបានណែនាំម៉ូឌុល \textenglish{EMAc2f} ថ្មី ដែលធ្វើសមាហរណកម្ម \textenglish{Efficient Multi-Scale Attention (EMA)} សម្រាប់ការកែលម្អលក្ខណៈពិសេសដោយដឹងពីភាពបាំង ជាមួយនឹងប្រតិបត្តិការ \textenglish{GhostBottleneck} សម្រាប់ការប្រើប្រាស់លក្ខណៈពិសេសឡើងវិញដោយចំណេញប៉ារ៉ាម៉ែត្រ។ លើសពីនេះ ការបង្កើនគុណភាពរូបភាពឌីណាមិក \textenglish{DySample} ត្រូវបានបញ្ចូលយ៉ាងយុទ្ធសាស្ត្រទៅក្នុងផ្នែក \textenglish{Neck} នៃបណ្តាញ ដើម្បីរក្សាព័ត៌មានលម្អិតនៃលំហប្រេកង់ខ្ពស់អំឡុងពេលផ្សារភ្ជាប់លក្ខណៈពិសេស ដែលជាការកាត់បន្ថយភាពព្រាលនៃវាយនភាពលម្អិតនៃលោហៈដោយផ្ទាល់។\\

ម៉ូដែលដែលត្រូវបានបង្កើនប្រសិទ្ធភាពនេះ ត្រូវបានដាក់ឱ្យដំណើរការលើ \textenglish{NVIDIA Jetson Orin Nano} ដោយប្រើការបង្កើនល្បឿន \textenglish{TensorRT} ភ្ជាប់ជាមួយ \textenglish{ByteTrack} សម្រាប់ភាពស៊ីសង្វាក់គ្នានៃពេលវេលា និងតក្កវិជ្ជានៃការរាប់តាមធរណីមាត្រ ដើម្បីការពារការរាប់ស្ទួន។ លទ្ធផលពិសោធន៍បានបង្ហាញថា \textenglish{YOLOv8-GED} សម្រេចបាននូវដំណើរការទាន់ពេលវេលាដ៏អស្ចារ្យ ដោយរក្សាបាន $89.9–100.8 FPS$ ជាមួយនឹងរយៈពេលពន្យារ (Latency) តិចជាង $11.5 ms$ និងការប្រើប្រាស់ផ្នែករឹងទាប $(GPU \approx 37–40 \%, CPU \approx 8 \%)$។ ចំណុចសំខាន់គឺថា វារក្សាបាននូវភាពជាក់លាក់ខ្ពស់នៃការរាប់ នៅក្នុងស្ថានភាពដែលមានដង់ស៊ីតេខ្ពស់ និងមានការបាំង ដែលកន្លែងម៉ូដែលគោលដែលមានទម្ងន់ធ្ងន់ជាងដូចជា \textenglish{RT-DETR} មានការធ្លាក់ចុះយ៉ាងខ្លាំង។\\

ក្រៅពីការអនុវត្តក្នុងឧស្សាហកម្មដោយផ្ទាល់ ការងារនេះបានរួមចំណែកនូវវិធីសាស្ត្រដែលអាចអនុវត្តឡើងវិញបានពីការរចនាទៅការដាក់ឱ្យដំណើរការ សម្រាប់ការចាប់សញ្ញាវត្ថុខ្នាតតូច នៅក្រោមកម្រិតធនធានគណនា Edge ដ៏តឹងរ៉ឹង។ ដោយត្រូវបានផ្ទៀងផ្ទាត់តាមរយៈ \textenglish{Proof of Concept (PoC)} ជាក់ស្តែងដែលទទួលបានជោគជ័យ វាផ្តល់នូវភស្តុតាងតាមរយៈការពិសោធន៍ថា ការកែប្រែស្ថាបត្យកម្មជាក់លាក់តាមវិស័យដែលមានគោលដៅច្បាស់លាស់ អាចធ្វើឱ្យប្រសើរឡើងយ៉ាងខ្លាំងនូវកម្រិត \textenglish{Recall} និងភាពរឹងមាំនៃការចាប់សញ្ញាវត្ថុខ្នាតតូច ដោយមិនលះបង់ល្បឿនដំណើរការទាន់ពេលវេលាឡើយ។
\end{khmer}
